\section*{Chapitre \ref{ch:b2:diffeq} --- Équations différentielles}

\begin{solution}{exo:b2:diffeq:1}
Première matrice~: \omterm{def:b2:reduction:eigen}{valeurs propres} $1, 2$, vecteurs propres $(1,0)$
et $(1,1)$. Décomposons $X_0 = (1,0) = 1\cdot(1,0) + 0\cdot(1,1)$~:
la solution est
\[
X(t) = \eu^{t}\begin{pmatrix}1\\ 0\end{pmatrix} .
\]
(Le vecteur initial est lui-même un vecteur propre.)

Seconde~: $A = 2I + N$, $N = E_{12}$, $N^2 = 0$~: $\eu^{tA} =
\eu^{2t}(I + tN)$, donc
\[
X(t) = \eu^{2t}\begin{pmatrix} 1 & t\\ 0 & 1\end{pmatrix}
\begin{pmatrix}1\\ 0\end{pmatrix}
= \eu^{2t}\begin{pmatrix}1\\ 0\end{pmatrix}.
\]
\end{solution}

\begin{solution}{exo:b2:diffeq:2}
$y' = \sin(ty)$~: $\abs{\sin(ty) - \sin(tz)} \leq \abs t\,\abs{y -
z}$ --- \omterm{def:b2:metric:continuity}{lipschitzienne} en $y$ \omterm{def:b2:funcseq:def}{uniformément} sur tout segment de
temps~: solutions globales uniques sur $\R$ (appliquer le théorème
sur chaque segment).

$y' = y^2$~: seulement localement \omterm{def:b2:metric:continuity}{lipschitzienne}~; pas de théorème
global, et en effet $y(0) = 1$ explose en $t = 1$.

$y' = \abs y$~: $\abs\cdot$ est $1$-lipschitzienne~: existence et
unicité globales. Avec $y(0) = 0$~: $y \equiv 0$ (unicité~!). Avec
$y(0) = 1$~: $y$ reste positive (ne peut traverser la solution
nulle), donc $y' = y$~: $y = \eu^t$.
\end{solution}

\begin{solution}{exo:b2:diffeq:3}
Si $y(t_1) = z(t_1)$ à un certain instant, alors $y$ et $z$
résolvent le même problème de Cauchy en $t_1$~: l'unicité impose $y
= z$ sur leur intervalle commun --- des solutions distinctes ne se
rencontrent jamais.

Pour $y' = y(1-y)$~: les constantes $0$ et $1$ sont solutions. Une
solution partant dans $\intoo{0}{1}$ ne peut jamais atteindre $0$
ni $1$ (elle heurterait une solution constante)~: elle reste dans
$\intoo{0}{1}$, donc est globale (bornée~: pas d'explosion --- par
exemple par le critère de l'\cref{exo:b2:diffeq:9}, ou parce que le
champ de vecteurs est borné sur la bande piège).
\end{solution}

\begin{solution}{exo:b2:diffeq:4}
$(A - 2I)^2 = \begin{pmatrix}1 & 1\\ -1 & -1\end{pmatrix}^2 = 0$~:
Dunford avec $D = 2I$, $N = A - 2I$~:
\[
\eu^{tA} = \eu^{2t}\,(I + tN)
= \eu^{2t}\begin{pmatrix} 1 + t & t\\ -t & 1 - t\end{pmatrix}.
\]
Variation des constantes avec $B(s) = (\eu^{2s}, 0)^{\mathsf T}$~:
\[
X(t) = \int_0^t \eu^{(t-s)A}B(s)\,\dd s
= \eu^{2t}\int_0^t \begin{pmatrix} 1 + (t-s)\\ -(t-s)
\end{pmatrix}\dd s
= \eu^{2t}\begin{pmatrix} t + \frac{t^2}{2}\\[2pt] -\frac{t^2}{2}
\end{pmatrix},
\]
en utilisant $\eu^{(t-s)A}B(s) = \eu^{2(t-s)}(I + (t-s)N)\,(\eu^{2s},
0)^{\mathsf T} = \eu^{2t}\bigl(1 + (t-s),\, -(t-s)\bigr)^{\mathsf
T}$. (Vérification~: $X(0) = 0$~; $X' - AX = (\eu^{2t}, 0)^{\mathsf
T}$ par dérivation.)
\end{solution}

\begin{solution}{exo:b2:diffeq:5}
\emph{Série~:} $A = \alpha I + \beta J$ avec $J =
\begin{pmatrix}0 & -1\\ 1 & 0\end{pmatrix}$, $J^2 = -I$~; les deux
termes commutent, donc $\eu^{tA} = \eu^{\alpha t}\,\eu^{\beta tJ}$,
et la série de $\eu^{\beta t J}$ se scinde selon les puissances
paires/impaires en $\cos(\beta t)I + \sin(\beta t)J$~: la matrice de
rotation-homothétie annoncée.

\emph{Complexe~:} identifions $(x, y) \in \R^2$ avec $z = x + \iu
y$~; le système $X' = AX$ s'écrit $z' = (\alpha + \iu\beta)z$, dont
la solution $z(t) = \eu^{\alpha t}\eu^{\iu\beta t}z_0$ est
exactement la spirale~: module $\eu^{\alpha t}$, argument avançant à
la vitesse $\beta$.
\end{solution}

\begin{solution}{exo:b2:diffeq:6}
Posons $v(t) = \eu^{-k(t - t_0)}\int_{t_0}^t u$. Alors
\[
v'(t) = \eu^{-k(t-t_0)}\Bigl(u(t) - k\int_{t_0}^t u\Bigr)
\leq C\,\eu^{-k(t-t_0)},
\]
par l'hypothèse. En intégrant de $t_0$ à $t$ ($v(t_0) = 0$)~:
$v(t) \leq \frac{C}{k}\bigl(1 - \eu^{-k(t - t_0)}\bigr)$,
c'est-à-dire $k\int_{t_0}^t u \leq C\bigl(\eu^{k(t-t_0)} - 1\bigr)$~;
en réinjectant dans l'hypothèse~: $u(t) \leq C\eu^{k(t-t_0)}$.

Unicité/dépendance~: deux solutions $y, z$ de l'équation intégrale
vérifient
\[
\norm{y(t) - z(t)} \leq \norm{y_0 - z_0} +
k\int_{t_0}^{t}\norm{y - z},
\]
et Gronwall avec $C = \norm{y_0 - z_0}$ donne la borne
exponentielle~; $C = 0$ donne l'unicité.
\end{solution}

\begin{solution}{exo:b2:diffeq:7}
Substituons $y = y_1 z$ avec $y_1 = \frac{\sin t}{t}$~: la formule
générale d'abaissement (\cref{prop:b2:diffeq:secondorder}) donne,
pour $u = z'$,
\[
y_1 u' + \Bigl(2y_1' + \frac{2}{t}\,y_1\Bigr)u = 0
\]
(l'équation normalisée en $y'' + \frac2t y' + y = 0$). Calculons
$2y_1' + \frac2t y_1 = 2\,\frac{t\cos t - \sin t}{t^2} +
\frac{2\sin t}{t^2} = \frac{2\cos t}{t}$~: ainsi
\[
\frac{u'}{u} = -\frac{2\cos t}{t}\cdot\frac{t}{\sin t}
= -2\cot t
\quad\Longrightarrow\quad
u = \frac{1}{\sin^2 t}
\quad (\text{à une constante près}),
\]
et $z = -\cot t$, d'où $y_2 = y_1 z = -\frac{\cos t}{t}$.
Solution générale sur $\intoo{0}{\pi}$~:
\[
y(t) = \alpha\,\frac{\sin t}{t} + \beta\,\frac{\cos t}{t} .
\]
(Ce sont les fonctions de Bessel sphériques d'\omterm{def:b2:structures:generated}{ordre} zéro.)
\end{solution}

\begin{solution}{exo:b2:diffeq:8}
Dunford~: $A = D + N$ qui commutent, $D$ \omterm{def:b2:reduction:diag}{diagonalisable} de mêmes
\omterm{def:b2:reduction:eigen}{valeurs propres}, $N$ nilpotent, donc
\[
\eu^{tA} = \eu^{tD}\,\eu^{tN},
\qquad
\eu^{tN} = \sum_{k < n} \frac{t^kN^k}{k!}
\ \text{(un polynôme matriciel en } t).
\]
Posons $-2\alpha = \max_i \Re\lambda_i < 0$. Dans une base
diagonalisant $D$, $\vertiii{\eu^{tD}} \leq \eu^{-2\alpha t}$
(coefficients $\eu^{t\lambda_i}$ de module $\eu^{t\Re\lambda_i}$)~;
les \omterm{def:b2:nvs:norm}{normes} dans des bases différentes diffèrent d'une constante.
D'où
\[
\norm{X(t)} \leq \vertiii{\eu^{tA}}\,\norm{X_0}
\leq C'\,\eu^{-2\alpha t}\,(1 + t)^{n-1}\,\norm{X_0}
\leq C\,\eu^{-\alpha t}\norm{X_0} ,
\]
en absorbant le polynôme dans un unique facteur exponentiel
($\eu^{-\alpha t}(1+t)^{n-1} \to 0$, donc borné).
\end{solution}

\begin{solution}{exo:b2:diffeq:9}
Soit $y$ une solution maximale sur $\intco{0}{T}$, $T \leq \infty$,
et supposons $T < \infty$. La forme intégrale donne, pour $t < T$,
\[
\norm{y(t)} \leq \norm{y_0} + \int_0^t \bigl(a\norm{y(s)} +
b\bigr)\dd s
\leq \bigl(\norm{y_0} + bT\bigr) + a\int_0^t\norm{y(s)}\,\dd s ,
\]
et Gronwall borne $\norm{y(t)} \leq (\norm{y_0} +
bT)\,\eu^{aT} =: M$ sur $\intco{0}{T}$~: la solution reste dans une
boule compacte. Alors $y' = f(t, y)$ est bornée au voisinage de
$T$, donc $y$ est \omterm{def:b2:metric:continuity}{lipschitzienne} au voisinage de $T$ et se prolonge
continûment en $T$ (critère de Cauchy)~; résoudre le problème de
Cauchy en $(T, y(T))$ prolonge $y$ au-delà de $T$, contredisant la
maximalité. Donc $T = \infty$~: pas d'échappement en temps fini
sous croissance linéaire.
\end{solution}

\begin{solution}{exo:b2:diffeq:10}
Racines caractéristiques de $r^2 - 3r + 2$~: $1$ et $2$, donc les
solutions homogènes sont $a\eu^t + b\eu^{2t}$. La supposition
$\alpha\eu^t$ échoue car $\eu^t$ résout déjà l'équation homogène
(la racine $r = 1$ «~entre en \omterm{pb:b2:diffeq:1}{résonance}~» avec le second membre).
Avec $y = \alpha t\,\eu^t$~: $y' = \alpha(1 + t)\eu^t$, $y'' =
\alpha(2 + t)\eu^t$, et
\[
y'' - 3y' + 2y = \alpha\eu^t\bigl(2 + t - 3 - 3t + 2t\bigr)
= -\alpha\,\eu^t :
\]
$\alpha = -1$, $y_p = -t\,\eu^t$. Solution générale~: $y = a\eu^t
+ b\eu^{2t} - t\eu^t$. Données initiales $y(0) = y'(0) = 0$~: $a + b
= 0$ et $a + 2b - 1 = 0$~: $b = 1$, $a = -1$~:
\[
y(t) = \eu^{2t} - (1 + t)\,\eu^{t} .
\]

\end{solution}

\begin{solution}{exo:b2:diffeq:11}
$A = \lambda I + N$ avec $N = E_{12} + E_{23}$~: $N^2 = E_{13}$,
$N^3 = 0$, et $\lambda I$ commute avec $N$~:
\[
\eu^{tA} = \eu^{\lambda t}\Bigl(I + tN +
\frac{t^2}{2}N^2\Bigr)
= \eu^{\lambda t}\begin{pmatrix}
1 & t & \frac{t^2}{2}\\
0 & 1 & t\\
0 & 0 & 1
\end{pmatrix}.
\]
Solutions~: $X(t) = \eu^{\lambda t}\bigl(X_0 + tNX_0 +
\frac{t^2}2N^2X_0\bigr)$ --- chaque composante est $\eu^{\lambda
t}$ fois un polynôme de degré $\leq 2$. La borne sur le degré est
l'indice de nilpotence moins un~: la série de $\eu^{tN}$ se tronque
à $N^2$.
\end{solution}

\begin{solution}{exo:b2:diffeq:12}
\begin{enumerate}
  \item Si $X(T) = X(0)$, alors $Y(t) = X(t + T)$ résout $Y' =
        AY + B(t + T) = AY + B(t)$ avec $Y(0) = X(0)$~:
        l'unicité (\cref{thm:b2:diffeq:linear}) donne $Y = X$,
        c'est-à-dire $X$ est $T$-périodique. La réciproque est
        triviale.
  \item Trigonaliser sur $\C$~: $A = PT'P^{-1}$ avec $T'$
        triangulaire supérieure, de diagonale $(\lambda_i)$.
        Toute puissance d'une matrice triangulaire est
        triangulaire de diagonale $(\lambda_i^k)$, donc $\eu^{TA}
        = P\eu^{TT'}P^{-1}$ est triangulaire dans la même base,
        de diagonale $(\eu^{T\lambda_i})$~: ce sont les \omterm{def:b2:reduction:eigen}{valeurs
        propres}. Alors $I - \eu^{TA}$ est inversible ssi
        $\eu^{T\lambda} \neq 1$ pour toute \omterm{def:b2:reduction:eigen}{valeur propre},
        c'est-à-dire ssi $T\lambda \notin 2\iu\pi\Z$, ce qui est
        l'hypothèse énoncée.
  \item Variation des constantes~: $X(T) = \eu^{TA}X(0) +
        \int_0^T\eu^{(T-s)A}B(s)\dd s$, donc $X(T) = X(0)$
        s'écrit
        \[
        \bigl(I - \eu^{TA}\bigr)X(0) =
        \int_0^{T}\eu^{(T-s)A}B(s)\,\dd s ,
        \]
        qui a une unique solution $X(0)$ sous l'hypothèse
        d'inversibilité~: exactement une solution $T$-périodique.
        Pour l'oscillateur harmonique ($\lambda =
        \pm\iu\omega$), le cas exclu est $\omega T \in
        2\pi\Z$~: un forçage dont la période est un multiple de
        la période propre --- la \omterm{pb:b2:diffeq:1}{résonance}, comme le quantifie le
        problème du week-end.
\end{enumerate}
\end{solution}

\begin{solution}{pb:b2:diffeq:1}
\textbf{1.} Le \omterm{def:b2:reduction:charpoly}{polynôme caractéristique} est $\lambda^2 -
\tau\lambda + \delta$, de racines
$\frac{\tau\pm\sqrt\Delta}{2}$. Si $\delta < 0$ alors $\Delta =
\tau^2 - 4\delta > 0$ et les deux racines réelles ont pour produit
$\delta < 0$~: signes opposés. Si $\delta > 0$ et $\Delta \geq
0$~: racines réelles de produit $> 0$ et de somme $\tau$~: toutes
deux du signe de $\tau$. Si $\Delta < 0$~: paire conjuguée $\alpha
\pm \iu\beta$ avec $\alpha = \frac\tau2$, $\beta =
\frac{\sqrt{-\Delta}}2$.

\textbf{2.} $X(t) = a\,\eu^{\mu t}v_- + b\,\eu^{\lambda t}v_+$.
Les orbites avec $a = 0$ (resp.\ $b = 0$) courent le long de la
droite propre instable (resp.\ stable)~; toutes les autres ont
$\norm X \to \infty$ dans les deux sens du temps, asymptotes à $\R
v_+$ quand $t \to +\infty$ et à $\R v_-$ quand $t \to -\infty$~:
l'image de la selle. La bornitude sur $\R$ tout entier force $b = 0$
(sinon explosion en $+\infty$) et $a = 0$ (en $-\infty$)~: seulement
l'origine.

\textbf{3.} Avec $\mu < \lambda < 0$, les deux exponentielles
décroissent~: $X(t) \to 0$. Si $b \neq 0$, factorisons
$\eu^{\lambda t}$~:
\[
X(t) = \eu^{\lambda t}\bigl(b\,v_\lambda + a\,\eu^{(\mu -
\lambda)t}v_\mu\bigr),
\qquad \eu^{(\mu-\lambda)t} \to 0 :
\]
la direction de $X(t)$ tend vers $\R v_\lambda$, la direction
propre \emph{lente} --- toutes les orbites sauf l'axe rapide y
arrivent tangentiellement (le panneau de droite des portraits de
phase du chapitre).

\textbf{4.} Dans la base où $A = \begin{pmatrix} \alpha &
-\beta\\ \beta & \alpha\end{pmatrix}$
(\cref{exo:b2:diffeq:5}~; pour les systèmes de l'oscillateur de la
partie II, cette forme est atteinte par un changement de base réel
explicite), la solution est $\eu^{\alpha t}$ fois une rotation
d'angle $\beta t$~: des spirales logarithmiques, contractantes
lorsque $\alpha = \frac\tau2 < 0$, dilatantes lorsque $\tau > 0$,
et des courbes fermées (des ellipses dans les coordonnées d'origine)
lorsque $\tau = 0$~: le centre.

\textbf{5.} $\Delta = 0$ donne la \omterm{def:b2:reduction:eigen}{valeur propre} double $\lambda =
\frac\tau2$~; par Cayley--Hamilton
(\cref{thm:b2:reduction:cayleyhamilton}), $(A - \lambda I)^2 =
0$, donc $N = A - \lambda I$ est nilpotent, commute avec
$\lambda I$, et $\eu^{tA} = \eu^{\lambda t}(I + tN)$. Si $N =
0$~: $A = \lambda I$, toutes les demi-droites sont des orbites
(nœud en étoile). Si $N \neq 0$~: $X(t) = \eu^{\lambda t}(X_0 +
tNX_0)$, et pour $NX_0 \neq 0$ la direction \omterm{def:b2:integration:improper}{converge} vers l'unique
direction propre $\operatorname{im}N$~: nœud impropre. Ceci complète
l'image trace--déterminant.

\textbf{6.} $\tau = -2\zeta\omega$, $\delta = \omega^2 > 0$,
$\Delta = 4\omega^2(\zeta^2 - 1)$. Donc~: $0 < \zeta < 1$ donne
$\Delta < 0$, $\tau < 0$~: foyer stable~; $\zeta = 1$~: $\Delta =
0$~: nœud stable dégénéré~; $\zeta > 1$~: $\Delta > 0$, $\tau <
0$, $\delta > 0$~: nœud stable~; $\zeta = 0$~: $\tau = 0$,
$\delta > 0$~: centre. Un trajet vertical dans le plan à $\delta =
\omega^2$.

\textbf{7.} Racines $r = -\zeta\omega \pm
\omega\sqrt{\zeta^2-1}$. Pour $\zeta < 1$~: $r = -\zeta\omega
\pm \iu\omega_d$, $\omega_d = \omega\sqrt{1-\zeta^2}$~:
\[
x(t) = \eu^{-\zeta\omega t}\bigl(a\cos\omega_dt +
b\sin\omega_dt\bigr)
= R\,\eu^{-\zeta\omega t}\cos(\omega_dt - \varphi) .
\]
Pour $\zeta = 1$~: $x = (a + bt)\eu^{-\omega t}$. Pour $\zeta >
1$~: $x = a\eu^{r_-t} + b\eu^{r_+t}$, les deux taux négatifs. Les
maxima successifs de $\abs x$ dans le cas sous-amorti se produisent
à des instants séparés de la pseudo-période $\frac{2\pi}{\omega_d}$
(même phase du cosinus), et leur rapport est
$\eu^{-\zeta\omega\cdot2\pi/\omega_d} =
\eu^{-2\pi\zeta/\sqrt{1-\zeta^2}}$~: le décrément logarithmique,
un indicateur d'amortissement lisible sur un oscilloscope.

\textbf{8.} En rationalisant,
\[
\abs{\lambda_{\mathrm{slow}}} = \omega\bigl(\zeta -
\sqrt{\zeta^2-1}\bigr)
= \frac{\omega}{\zeta + \sqrt{\zeta^2 - 1}} ,
\]
dont le dénominateur croît avec $\zeta \geq 1$~: le taux de
décroissance est maximal en $\zeta = 1$, où il vaut $\omega$. Une
porte sur-amortie se ferme sans claquer mais \emph{lentement}~;
l'amortissement critique est l'optimum de l'ingénieur.

\textbf{9.} $E' = x'x'' + \omega^2xx' =
x'\bigl(-2\zeta\omega x' - \omega^2x\bigr) + \omega^2xx' =
-2\zeta\omega\,x'^2 \leq 0$. Si $x$ était périodique et
non constante, $E$ serait périodique et décroissante, donc
constante, forçant $x' \equiv 0$~: $x$ constante, et alors
$\omega^2x = 0$~: $x \equiv 0$. Donc pour $\zeta > 0$ la seule
solution périodique est le repos~: l'amortissement tue tout \omterm{def:b2:structures:sn}{cycle}.

\textbf{10.} Le centre $\zeta = 0$ vit sur la droite $\tau =
0$ du plan trace--déterminant --- un ensemble d'intérieur vide~:
une perturbation arbitrairement petite de la matrice (tout
amortissement physique) déplace $\tau$ hors de zéro et transforme
les orbites fermées en spirales. La périodicité de l'oscillateur
non amorti est donc un phénomène sur le fil du rasoir, non un
phénomène robuste.

\textbf{11.} En substituant $x_p = \Re(z\eu^{\iu\gamma t})$ dans
l'équation~:
\[
\bigl(-\gamma^2 + 2\iu\zeta\omega\gamma +
\omega^2\bigr)z = F
\quad\Longrightarrow\quad
z = \frac{F}{\omega^2 - \gamma^2 + 2\iu\zeta\omega\gamma},
\]
donc $x_p = \abs z\cos(\gamma t - \varphi)$ avec $\varphi =
\arg(\omega^2 - \gamma^2 + 2\iu\zeta\omega\gamma)$, c'est-à-dire
$\tan\varphi = \frac{2\zeta\omega\gamma}{\omega^2 -
\gamma^2}$, et le $A(\gamma) = \abs z$ annoncé.

\textbf{12.} La différence de deux solutions résout l'équation
homogène, qui pour $\zeta > 0$ décroît vers $0$
(question 7)~: toute solution est égale à $x_p$ plus un transitoire
s'évanouissant à l'infini. Le régime permanent est un attracteur
global~: les conditions initiales sont oubliées, seuls subsistent
$A(\gamma)$ et le déphasage $\varphi$.

\textbf{13.} Minimisons $g(u) = (\omega^2 - u)^2 +
4\zeta^2\omega^2u$ sur $u = \gamma^2 \geq 0$~: $g'(u) =
-2(\omega^2 - u) + 4\zeta^2\omega^2 = 0$ en $u = \omega^2(1 -
2\zeta^2)$, intérieur ssi $\zeta < \frac{1}{\sqrt2}$. Là,
\[
g(u_*) = 4\zeta^4\omega^4 + 4\zeta^2\omega^4(1 - 2\zeta^2)
= 4\zeta^2\omega^4(1 - \zeta^2),
\qquad
A(\gamma_*) = \frac{F}{2\zeta\omega^2\sqrt{1 - \zeta^2}} .
\]
Face à la réponse statique $A(0) = \frac{F}{\omega^2}$~:
amplification $\frac{1}{2\zeta\sqrt{1-\zeta^2}} \approx
\frac{1}{2\zeta}$ pour $\zeta$ petit --- un système faiblement
amorti proche de $\gamma_* \approx \omega$ multiplie l'entrée au
centuple lorsque $\zeta = 0.005$.

\textbf{14.} Le $x$ annoncé vérifie $x(0) = x'(0) = 0$ et
\[
x'' + \omega^2 x = \frac{F(\omega^2 -
\gamma^2)\cos\gamma t}{\omega^2 - \gamma^2} = F\cos\gamma t
\]
(les parties en $\cos\omega t$ se compensent). La forme produit
découle de $\cos p - \cos q = 2\sin\frac{q+p}{2}\sin\frac{q-p}2$
avec $p = \gamma t$, $q = \omega t$. Pour $\gamma$ proche de
$\omega$, le facteur $\sin\frac{(\omega-\gamma)t}2$ est une
enveloppe lente modulant l'oscillation rapide
$\sin\frac{(\omega+\gamma)t}2$~: des battements, d'amplitude
$\frac{2F}{\abs{\omega^2-\gamma^2}}$ --- grande, mais bornée.

\textbf{15.} Pour $x_p = \frac{F}{2\omega}t\sin\omega t$~:
\[
x_p'' = \frac{F}{2\omega}\bigl(2\omega\cos\omega t -
\omega^2t\sin\omega t\bigr)
= F\cos\omega t - \omega^2x_p :
\]
une solution. Et à $t$ fixé, en faisant $\gamma \to \omega$ dans
la question 14~:
\[
\frac{2F\sin\frac{(\omega-\gamma)t}2
\sin\frac{(\omega+\gamma)t}{2}}
{(\omega-\gamma)(\omega+\gamma)}
\longrightarrow
\frac{2F\cdot\frac{(\omega-\gamma)t}2\big/(\omega-\gamma)
\cdot\sin\omega t}{2\omega}
= \frac{F\,t\sin\omega t}{2\omega} .
\]
L'amplitude croît linéairement sans borne~: la catastrophe de
\omterm{pb:b2:diffeq:1}{résonance} --- la raison pour laquelle les soldats rompent le pas
sur les ponts.

\textbf{16.} Le signal carré porte des harmoniques à toute
fréquence impaire $n = 1, 3, 5, \dots$~; par linéarité, chaque
harmonique $n$ est amplifiée par la réponse de l'oscillateur en
$\gamma = n$. Pour $\omega = 3$, la troisième harmonique atteint
exactement la \omterm{pb:b2:diffeq:1}{résonance}. Les ingénieurs craignent les entrées
carrées (et en dents de scie) parce qu'elles excitent \emph{toutes}
les harmoniques impaires à la fois~: quelle que soit la fréquence
propre de la structure, une harmonique l'attend.

\textbf{17.} $W' = y_1y_2'' - y_1''y_2 = -qy_1y_2 + qy_1y_2 =
0$~: $W$ est constant (Liouville avec une matrice compagnon de
trace nulle). $W = 0$ en un point rend les données initiales de
$y_2$ proportionnelles à celles de $y_1$, donc $y_2$ proportionnelle
à $y_1$ (unicité)~; $W \neq 0$ ssi indépendantes. Si $y(t_0) =
y'(t_0) = 0$ alors $y \equiv 0$ (unicité)~: une solution non nulle a
des zéros simples, et un zéro simple est isolé ($y'$ de signe fixe
au voisinage).

\textbf{18.} Entre deux zéros consécutifs $a < b$, $y_1$ garde
un signe constant, disons $y_1 > 0$ sur $\intoo ab$~: alors
$y_1'(a) > 0$ et $y_1'(b) < 0$ (zéros simples). En évaluant la
constante $W = y_1y_2' - y_1'y_2$ en $a$ et $b$~:
\[
W = -y_1'(a)\,y_2(a) = -y_1'(b)\,y_2(b) ,
\]
donc $y_2(a)$ et $y_2(b)$ sont de signes opposés ($W \neq 0$
interdit à l'un ou l'autre de s'annuler)~: $y_2$ s'annule dans
$\intoo ab$ (valeurs intermédiaires). Il ne peut s'y annuler deux
fois~: deux zéros de $y_2$ encadreraient un zéro de $y_1$ par le
même argument avec les rôles échangés, contredisant la
consécutivité~: exactement un zéro --- l'entrelacement.

\textbf{19.} Supposons que $z$ n'a aucun zéro dans $\intoo ab$~;
en remplaçant $y, z$ par leurs opposés, supposons $y > 0$ et $z >
0$ sur $\intoo ab$. Posons $\varphi = yz' - y'z$~: $\varphi' = yz''
- y''z = (q_1 - q_2)\,yz \leq 0$ sur $\intoo ab$~: $\varphi$ est
décroissante. Mais $\varphi(a) = -y'(a)z(a) \leq 0$ (car
$y'(a) > 0$, $z(a) \geq 0$) et $\varphi(b) = -y'(b)z(b) \geq
0$ (car $y'(b) < 0$, $z(b) \geq 0$)~: une fonction décroissante
allant de $\leq 0$ à $\geq 0$ s'annule identiquement, donc
$(q_1 - q_2)yz \equiv 0$ sur $\intoo ab$. Si $q_1 < q_2$
quelque part dans $\intoo ab$, c'est absurde ($y, z > 0$ en ce
point)~: $z$ doit s'annuler strictement à l'intérieur. En général
($q_1 \leq q_2$), ou bien $z$ s'annule dans $\intoo ab$, ou bien
$\varphi \equiv 0$ force $z$ proportionnelle à $y$, qui s'annule en
$a$ et $b$~: dans tous les cas $z$ a un zéro dans $\intcc ab$.

\textbf{20.} Borne supérieure~: comparons $y$ (coefficient $q \geq
m^2$) avec $u(t) = \sin(m(t - a))$ (coefficient $m^2 \leq q$,
donc $y$ joue le rôle de $z$ dans la question 19)~: si $y$ n'avait
aucun zéro dans $\intoc{a}{a + \pi/m}$, les zéros $a$ et $a +
\frac\pi m$ de $u$ seraient consécutifs avec $y \neq 0$
entre eux, contredisant la question 19~: deux zéros consécutifs de
$y$ sont à distance $\leq \frac\pi m$. Borne inférieure~: si deux
zéros consécutifs $a < b$ de $y$ vérifiaient $b - a < \frac\pi M$,
alors $z(t) = \sin(M(t-a))$ (coefficient $M^2 \geq q$) devrait
s'annuler dans $\intcc ab \subset \intoo{a}{a + \pi/M}
\cup\{a\}$, où son unique zéro est $a$ lui-même --- mais
la question 19 appliquée sur $\intoo{a}{b}$ avec stricte inégalité
aux extrémités donne un zéro dans $\intcc ab$, et $z > 0$ sur
$\intoc ab$~: contradiction. Donc $\frac\pi M \leq b - a \leq
\frac\pi m$~; pour $q \equiv \omega^2$ les deux bornes se rejoignent
à l'espacement exact $\frac\pi\omega$ de l'oscillateur
harmonique.

\textbf{21.} Avec $u = ty$~: $u'' = ty'' + 2y'$, donc $ty'' + 2y'
+ ty = u'' + u = 0$~: $u = A\sin t + B\cos t = R\sin(t +
\varphi)$, et $y = \frac{u}{t}$ retrouve $\frac{\sin t}t$ et
$\frac{\cos t}t$ (\cref{exo:b2:diffeq:7}) sans aucun abaissement de
l'\omterm{def:b2:structures:generated}{ordre}. Les zéros de toute solution non nulle sont ceux de
$R\sin(t + \varphi)$~: espacés exactement de $\pi$ --- la
philosophie de Sturm à l'œuvre~: le coefficient $q \equiv 1$ dicte
les zéros, avec ou sans formule.

\textbf{22.} Posons $x(t) =
\frac1\omega\int_0^t\sin(\omega(t-s))F(s)\dd s$. Alors $x(0) =
0$~;
\[
x'(t) = \frac1\omega\sin(0)F(t) +
\int_0^t\cos(\omega(t-s))F(s)\dd s
= \int_0^t\cos(\omega(t-s))F(s)\dd s ,
\]
donc $x'(0) = 0$~; et $x''(t) = F(t) -
\omega\int_0^t\sin(\omega(t-s))F(s)\dd s = F(t) - \omega^2x(t)$
(dérivation d'une \omterm{thm:b2:integration:continuity}{intégrale à paramètre} à borne variable, comme
dans le chapitre d'intégration). Avec $F(s) = F\cos\omega s$, la
transformation produit en somme donne
\[
\int_0^t\sin(\omega(t-s))\cos(\omega s)\dd s
= \frac12\int_0^t\bigl(\sin\omega t + \sin(\omega t -
2\omega s)\bigr)\dd s
= \frac{t}{2}\sin\omega t ,
\]
(le second morceau s'intègre en zéro), donc $x =
\frac{F}{2\omega}t\sin\omega t$~: de nouveau la question 15, à
partir de Duhamel.

\textbf{23.} Sous forme de système $X' = AX + (0, F(t))^{\mathsf T}$
avec $\operatorname{Sp}A$ de parties réelles négatives ($\zeta >
0$)~: la variation des constantes et l'\cref{exo:b2:diffeq:8}
($\vertiii{\eu^{tA}} \leq C\eu^{-\alpha t}$) donnent
\[
\norm{X(t)} \leq C\eu^{-\alpha t}\norm{X_0}
+ \int_0^t C\eu^{-\alpha(t-s)}\norm{F}_\infty\dd s
\leq C\norm{X_0} + \frac{C\norm F_\infty}{\alpha} :
\]
entrée bornée, sortie bornée --- \omterm{def:b2:funcseq:def}{uniformément} en les données
initiales après le transitoire.

\textbf{24.} Pour $\zeta = 0$ et $\gamma \neq \omega$,
la solution de la question 14 est bornée, et ajouter une solution
homogène quelconque (bornée~: les orbites du centre sont des
cercles) la maintient bornée~; en $\gamma = \omega$, la question 15
croît linéairement. Ainsi, pour l'oscillateur non amorti, la
bornitude sous forçage périodique échoue à exactement une
fréquence --- la \omterm{pb:b2:diffeq:1}{résonance} --- tandis que la question 23 montre que
tout amortissement positif restaure la bornitude pour \emph{toutes}
les entrées bornées.

\textbf{25.} (i) Le plan trace--déterminant classe toute la
dynamique linéaire plane autonome, et l'oscillateur de la partie II
en parcourt une droite verticale~; le forçage quitte le plan
(non autonome), et Duhamel prend le relais. (ii) $\gamma_*$ est
la fréquence que le système préfère, $A(\gamma_*)$ le prix à payer
pour l'exciter, et $\frac{1}{2\zeta}$ le facteur d'amplification
--- l'acuité de la \omterm{pb:b2:diffeq:1}{résonance} que les ingénieurs appellent le
facteur de qualité. (iii) Les théorèmes de Sturm lisent
l'oscillation sur le seul signe et la seule taille de $q$~: ils
gouvernent des équations (Bessel, Schrödinger) dont les solutions
n'ont pas de formule élémentaire. (iv) Les \omterm{def:b2:diffeq:wronskian}{wronskiens} constants
(question 17, via Liouville) et la stabilité entrée bornée
(question 23) sont réutilisés discrètement chaque fois que le livre
rencontre des équations à coefficients variables ou des systèmes
perturbés.
\end{solution}
